\documentclass[10pt,twocolumn]{article}
\usepackage[T1]{fontenc}
\usepackage[utf8]{inputenc}
\usepackage{lmodern}
\usepackage[margin=1.8cm,top=2.4cm,bottom=2cm,columnsep=0.6cm]{geometry}
\usepackage{fancyhdr}
\usepackage{titlesec}
\usepackage{booktabs}
\usepackage{amsmath,amssymb,amsfonts}
\usepackage{microtype}
\usepackage{xcolor}
\usepackage{mdframed}
\usepackage{parskip}
\usepackage{enumitem}
\usepackage{hyperref}

\definecolor{rulered}{RGB}{180,30,30}
\definecolor{boxgray}{RGB}{245,245,245}
\definecolor{keywordgray}{RGB}{100,100,100}

\pagestyle{fancy}
\fancyhf{}
\fancyhead[L]{\textsc{\small Claude Mag}}
\fancyhead[C]{\small\textit{Machine Learning Research Digest}}
\fancyhead[R]{\small 2026-09-07}
\fancyfoot[C]{\small\thepage}
\renewcommand{\headrulewidth}{0.8pt}

\titleformat{\section}{\normalsize\bfseries\color{rulered}}{}{0pt}{}%
  [\vspace{-4pt}\color{rulered}\rule{\columnwidth}{0.4pt}]
\titlespacing{\section}{0pt}{8pt}{4pt}

\hypersetup{colorlinks=true,urlcolor=rulered,linkcolor=black}

\begin{document}
\setlength{\parindent}{0pt}
\setlength{\parskip}{4pt}

{\small\color{keywordgray}\textbf{Keywords:}
  vision-language-action models \textbullet{}
  embodied AI \textbullet{}
  robot manipulation \textbullet{}
  skill-based planning}\par\vspace{4pt}

{\LARGE\bfseries\raggedright
  EmbodiedSkills: Orchestrating, Training,
  and Deploying VLA Agents}\par\vspace{4pt}

{\large\itshape\raggedright
  A Unified Skill Interface with Precondition
  Checking and Postcondition Verification
  Lifts Robot Manipulation to 86\% Success}\par\vspace{6pt}

{\small\textbf{By} Wei Wang, Wenqiao Zhang, Yutong Lin, et al.\
  (17 authors; Zhejiang University et al.)}\par\vspace{2pt}

{\small\color{keywordgray}arXiv:2609.01281 \textbullet{} Preprint, September 2026}
\par\vspace{2pt}\noindent{\color{rulered}\rule{\columnwidth}{0.6pt}}\vspace{6pt}

Vision-Language-Action (VLA) models map visual observations and language
instructions to robot actions end-to-end, but they struggle with
long-horizon tasks because early errors cascade without a recovery
mechanism. EmbodiedSkills wraps VLA policies in an executable-skill
abstraction that checks prerequisites before execution and verifies
outcomes afterward. Task-adapted low-level VLA policies within this
framework achieve 86.20\% average success across 50 RoboTwin~2.0 tasks
and 97.40\% across four LIBERO suites --- far above frozen-policy baselines.

\begin{mdframed}[backgroundcolor=boxgray,linecolor=rulered,linewidth=1pt,
    innerleftmargin=6pt,innerrightmargin=6pt,
    innertopmargin=6pt,innerbottommargin=6pt]
\textbf{\small AT A GLANCE}\par\vspace{4pt}
\begin{description}[leftmargin=0pt,labelsep=4pt,itemsep=2pt,parsep=0pt,
    font=\small\bfseries]
  \item[Problem:] VLA models fail on long-horizon tasks because they
    have no mechanism to detect or recover from execution errors.
  \item[Why it matters:] Real-world robot deployment requires reliable
    multi-step manipulation; a single failure can ruin a long task.
  \item[Method:] Executable-skill interface with learned precondition
    and postcondition classifiers; modular VLA policy adaptation.
  \item[Result:] 86.20\% average success (RoboTwin 2.0), 97.40\%
    (LIBERO); 62\% reduction in cascading failures on 6+ skill tasks.
  \item[Evidence:] Comparison vs.\ frozen OpenVLA and task-adapted VLA
    without skill structure across 54 robot manipulation tasks.
\end{description}
\end{mdframed}

\section{The Big Picture}

VLA models have demonstrated impressive generalisation on short-horizon
manipulation tasks: given an image observation and a language description,
they generate motor commands that successfully complete simple pick-and-place
or push tasks. However, real-world applications require long-horizon
plans: setting a table, assembling a product, or restocking a shelf
involves sequences of 6--10 or more distinct sub-tasks.

In long-horizon settings, VLA models fail because they lack two critical
properties:
\begin{itemize}[itemsep=2pt,parsep=0pt]
  \item \textbf{Error detection:} the model has no way to know whether
    a sub-task succeeded or failed before moving to the next.
  \item \textbf{Recovery:} when a sub-task fails, the model cannot
    re-try or select an alternative action sequence.
\end{itemize}

EmbodiedSkills introduces both properties by structuring VLA execution
around \emph{executable skills} --- a design pattern borrowed from
classical task-and-motion planning but implemented entirely with learned
neural components.

\section{Why Existing Approaches Fall Short}

\textbf{Monolithic VLA policies} (e.g.\ OpenVLA, RT-2) are trained
end-to-end and have no explicit sub-task structure. They may hallucinate
task completion or proceed past failures.

\textbf{Hierarchical RL approaches} train separate high-level and
low-level policies, but the high-level policy must be retrained for
each new task, and low-level sub-policies typically cannot be reused
across hierarchies.

\textbf{Classical task planners (e.g.\ PDDL-based)} offer structured
execution with precondition/postcondition semantics but require
hand-coded predicates and cannot handle the visual complexity of
real-world robot environments.

EmbodiedSkills occupies the gap: learned precondition/postcondition
classifiers from visual observations, structured execution via a fixed
skill interface, and task-adapted VLA policies as low-level executors.

\section{Core Mechanism}

A skill $s$ is a tuple:
\[
  s = (\texttt{name},\; \texttt{pre}(o),\; \pi_s(o, g),\; \texttt{post}(o))
\]
where:
\begin{itemize}[itemsep=1pt,parsep=0pt]
  \item $\texttt{pre}(o) \to \{0,1\}$: checks whether the robot is in
    a valid state to begin executing $s$.
  \item $\pi_s(o, g) \to a$: the low-level VLA policy that generates
    actions to achieve sub-goal $g$ from observation $o$.
  \item $\texttt{post}(o) \to \{0,1\}$: checks whether skill $s$
    completed successfully.
\end{itemize}

The agent loop selects skills via a vision-language model planner, checks
preconditions, executes until postcondition is satisfied or a timeout is
reached, and triggers recovery or replanning on failure.

\section{Technical Formulation}

\subsection{High-Level Planning}

The high-level planner is a VLM queried at each step as:
\begin{equation}
  s_t, a_t = \mathrm{VLM}(G, o_t, \mathcal{S})
\end{equation}
where $G$ is the task goal, $o_t$ is the current observation, and
$\mathcal{S}$ is the skill library. The VLM reasons in natural language
via chain-of-thought before selecting the next skill $s_t$ and its
arguments $a_t$.

\subsection{Precondition and Postcondition Classifiers}

Pre- and postconditions are binary classifiers trained on annotated
demonstrations:
\begin{align}
  f_\mathrm{pre}(o;\,\theta_\mathrm{pre}) &\to P(\texttt{pre satisfied}) \\
  f_\mathrm{post}(o;\,\theta_\mathrm{post}) &\to P(\texttt{post satisfied})
\end{align}
with positive/negative labels derived from states before and after
successful vs.\ failed skill executions.

\subsection{Task-Adapted Policy Training}

Each skill's low-level policy $\pi_s$ is trained with a combination of
behavioural cloning and proximal policy optimisation:
\[
  \mathcal{L}_s(\theta) = \mathcal{L}_\mathrm{BC}(\theta)
  + \lambda\, \mathcal{L}_\mathrm{PPO}(\theta, R_s)
\]
where $R_s$ is a dense reward derived from the postcondition classifier:
\[
  R_s(o_t) = f_\mathrm{post}(o_t;\,\theta_\mathrm{post})
\]
This reward is dense --- it provides a training signal at every step
proportional to how close the current state is to satisfying the
postcondition --- avoiding the sparse reward problem common in robot RL.

\subsection{Recovery Policy}

Failure modes are classified by the postcondition checker; each failure
mode has an associated recovery policy $r$ selected from a predefined
recovery library:
\[
  r_t = \mathrm{RecoveryPolicy}(o_t, s_t, \mathrm{failure\_mode})
\]

\section{Learning or Inference Procedure}

EmbodiedSkills is trained in three stages:
\begin{enumerate}[itemsep=2pt,parsep=0pt]
  \item \textbf{Precondition/postcondition learning.} Classifiers are
    trained on annotated state transitions from expert demonstrations.
  \item \textbf{Low-level VLA adaptation.} Each skill policy is fine-tuned
    from a pretrained VLA (OpenVLA) using BC + PPO with the postcondition
    reward. Adaptation is independent per skill.
  \item \textbf{High-level planner.} A pretrained VLM (7B or 72B) is
    prompted with the skill library and task specification; no VLM
    fine-tuning is required.
\end{enumerate}

At inference, the agent loop runs: plan $\to$ check precondition $\to$
execute $\to$ verify postcondition $\to$ recover or continue.

\section{What the Guarantee Says}

The paper proves that under the learned classifier regime, the expected
task completion rate $P(\mathrm{task})$ satisfies:
\[
  P(\mathrm{task}) \geq \prod_{i=1}^{K} P(\mathrm{success}_i)
  \cdot (1 - P(\mathrm{cascade}))
\]
where $K$ is the number of skills, $P(\mathrm{success}_i)$ is the
per-skill success rate, and $P(\mathrm{cascade})$ is the probability of
a cascade failure that the recovery policy cannot resolve. Adding
postcondition verification suppresses $P(\mathrm{cascade})$ by catching
failures before the next skill begins.

\section{Experimental Findings}

\begin{center}
\begin{tabular}{lcc}
\toprule
Method & RoboTwin 2.0 & LIBERO \\
\midrule
Frozen VLA (OpenVLA)        & 51.4\% & 78.6\% \\
Task-adapted (no framework) & 71.2\% & 89.1\% \\
\textbf{EmbodiedSkills}     & \textbf{86.2\%} & \textbf{97.4\%} \\
\bottomrule
\end{tabular}
\end{center}

\textbf{Long-horizon tasks (8+ skills):} EmbodiedSkills achieves 74.3\%
completion vs.\ 31.2\% for the task-adapted baseline without skill
structure, a 2.4$\times$ improvement attributable to verification-driven
recovery.

\section{Ablations and Interpretation}

\textbf{Without postcondition verification:} Success drops from 86.2\%
to 61.4\% on RoboTwin~2.0 --- the largest single ablation contributor.
Cascading failures account for the majority of the 24.8-point gap.

\textbf{Without precondition checking:} Success drops to 74.8\%, showing
that preventing skill execution from an invalid start state avoids
a different class of compounding errors.

\textbf{VLM planner size:} Using a 7B planner vs.\ a 72B planner reduces
success by 4.1\% on RoboTwin~2.0. The structured skill interface
compensates for planner imperfection: the planner only needs to choose
the right skill, not the right low-level action.

\textbf{Fixed vs.\ adaptive skill library:} A fixed library of 28
primitive skills matches a dynamically grown library for the tested
tasks, suggesting that a compact skill vocabulary covers the space of
practical manipulation scenarios.

\textbf{Dense reward (postcondition) vs.\ sparse reward:} Using the
postcondition classifier as a dense reward (vs.\ sparse 0/1 completion
reward) increases average per-skill success rate by 11.4 points.

\section*{Reference}

Wei Wang, Wenqiao Zhang, Yutong Lin et al.\
\textbf{EmbodiedSkills: A Unified Framework for Orchestrating,
Training, and Deploying VLA Agents.}
arXiv:2609.01281, September 2026.
\url{https://arxiv.org/abs/2609.01281}

\end{document}
